%% ============================================================
%%  Informe Proyecto Final — IA · EAFIT 2026-1
%%  Sistema de Predicción de Abandono Estudiantil
%%  Compilar con pdfLaTeX en Overleaf
%% ============================================================

\documentclass[10pt,twocolumn]{article}

\usepackage[utf8]{inputenc}
\usepackage[T1]{fontenc}
\usepackage[spanish]{babel}
\usepackage[margin=1.8cm, top=2.2cm, bottom=2.2cm]{geometry}
\usepackage{lmodern}
\usepackage{microtype}
\usepackage{palatino}
\usepackage{amsmath, amssymb}
\usepackage{graphicx}
\usepackage[dvipsnames,svgnames]{xcolor}
\usepackage{tikz}
\usepackage{pgfplots}
\pgfplotsset{compat=1.18}
\usepackage{booktabs}
\usepackage{tabularx}
\usepackage{multirow}
\usepackage{array}
\usepackage{listings}
\usepackage{algorithm}
\usepackage{algpseudocode}
\usepackage{hyperref}
\usepackage[numbers,sort&compress]{natbib}
\usepackage{enumitem}
\usepackage{float}
\usepackage{caption}
\usepackage{subcaption}
\usepackage{fancyhdr}
\usepackage{titlesec}

%% -- Paleta EAFIT -------------------------------------------
\definecolor{eafitblue}{RGB}{0, 61, 121}
\definecolor{eafitgreen}{RGB}{0, 150, 80}
\definecolor{eafitgray}{RGB}{90, 90, 90}
\definecolor{codebg}{RGB}{248, 248, 252}
\definecolor{codecomment}{RGB}{100, 110, 130}
\definecolor{codekw}{RGB}{0, 90, 170}
\definecolor{codestr}{RGB}{160, 50, 0}

\lstset{
  backgroundcolor=\color{codebg},
  basicstyle=\ttfamily\scriptsize,
  keywordstyle=\color{codekw}\bfseries,
  commentstyle=\color{codecomment}\itshape,
  stringstyle=\color{codestr},
  numbers=left, numberstyle=\tiny\color{eafitgray},
  numbersep=6pt, frame=single, framesep=3pt,
  rulecolor=\color{eafitblue!30},
  breaklines=true, captionpos=b, language=Python,
  showstringspaces=false, tabsize=4,
}

\titleformat{\section}
  {\normalfont\bfseries\large\color{eafitblue}}
  {\thesection.}{0.6em}{}
  [\vspace{-2pt}\color{eafitblue!35}\hrule\vspace{4pt}]

\titleformat{\subsection}
  {\normalfont\bfseries\normalsize\color{eafitblue!80}}
  {\thesubsection.}{0.5em}{}

\pagestyle{fancy}
\fancyhf{}
\fancyhead[L]{\small\color{eafitgray}\textit{Proyecto Final . Inteligencia Artificial . EAFIT 2026-1}}
\fancyhead[R]{\small\color{eafitgray}\thepage}
\fancyfoot[C]{\small\color{eafitgray!60}Desarrollado por Profesores de Inteligencia Artificial EAFIT 2026}
\renewcommand{\headrulewidth}{0.4pt}
\renewcommand{\headrule}{%
  \hbox to\headwidth{\color{eafitblue!35}\leaders\hrule height \headrulewidth\hfill}}

\hypersetup{
  colorlinks=true, linkcolor=eafitblue,
  citecolor=eafitgreen, urlcolor=eafitblue!80,
  pdftitle={Sistema de Prediccion de Abandono Estudiantil — EAFIT 2026},
  pdfauthor={Equipo — EAFIT},
}

\captionsetup{font=small, labelfont={bf,color=eafitblue}, skip=4pt}

\newcommand{\nota}[1]{%
  \vspace{4pt}
  \noindent\colorbox{eafitblue!8}{%
    \begin{minipage}{\columnwidth-2\fboxsep}
      \small\color{eafitblue!80}\textbf{Nota:} #1
    \end{minipage}
  }
  \vspace{4pt}
}

% ============================================================
\begin{document}

%% ==========================================================
%%  PORTADA
%% ==========================================================
\twocolumn[{%
\begin{center}
  {\color{eafitblue!25}\rule{\linewidth}{3pt}}\\[6pt]

  {\Large\bfseries\color{eafitblue}
    Sistema de Predicción de Abandono Estudiantil\\[4pt]
    mediante Machine Learning y Recomendaciones Personalizadas\\[4pt]
  }

  {\normalsize\color{eafitgray}
    Clasificación binaria con Random Forest y generación de recomendaciones
    para intervención temprana\\[10pt]
  }

  \begin{tabular}{cc}
    \textbf{Santiago Arboleda} & \textbf{Juan Esteban Villada} \\
    \small\texttt{sarboledag@eafit.edu.co} & \small\texttt{Jevilladao@eafit.edu.co}\\[6pt]
    \multicolumn{2}{c}{\textbf{Cristian Cabezas}} \\
    \multicolumn{2}{c}{\small\texttt{cdcabezasj@eafit.edu.co}}
  \end{tabular}\\[10pt]

  {\small\color{eafitgray}
    Escuela de Ciencias Aplicadas e Ingeniería . Universidad EAFIT\\
    Curso: Inteligencia Artificial . Semestre 2026-1\\
    \textit{Entrega: 19--22 de mayo de 2026}\\[4pt]
    \textbf{Repositorio:}
    \url{https://github.com/sarboledag/PROYECTO-FINAL}
  }\\[6pt]
  {\color{eafitblue!25}\rule{\linewidth}{1pt}}\\[10pt]
\end{center}
}]

%% ==========================================================
%%  ABSTRACT
%% ==========================================================
\begin{abstract}
\noindent
El abandono estudiantil representa un problema relevante para las
instituciones de educación superior, ya que afecta la permanencia académica,
la eficiencia institucional y el acompañamiento oportuno de los estudiantes.
Este proyecto desarrolla un sistema de Inteligencia Artificial para predecir
el riesgo de abandono estudiantil a partir de variables académicas, económicas
y personales.

Se utilizó el dataset \textit{Predict Students' Dropout and Academic Success},
transformando la variable objetivo original en una tarea de clasificación
binaria: abandono y no abandono. Se realizó un análisis exploratorio de datos,
preprocesamiento, entrenamiento de modelos y evaluación con métricas de
clasificación. Se compararon tres enfoques: un baseline basado en la clase
mayoritaria, una Regresión Logística y un Random Forest.

El modelo Random Forest fue seleccionado como modelo principal, alcanzando
un Accuracy de 0.8847, un Recall de 0.7394, un F1-score de 0.8046 y un
AUC-ROC de 0.9316. Además, se implementó una aplicación en Streamlit que
permite seleccionar un estudiante, estimar su probabilidad de abandono y
generar una recomendación personalizada para intervención temprana.\\[4pt]

\textbf{Palabras clave:} abandono estudiantil, Machine Learning,
Random Forest, clasificación binaria, Streamlit, recomendación personalizada.
\end{abstract}

\vspace{4pt}

%% ==========================================================
%%  1. INTRODUCCIÓN
%% ==========================================================
\section{Introducción}

El abandono estudiantil es un fenómeno que afecta a muchas instituciones
de educación superior. Cuando un estudiante abandona su proceso académico,
no solo se ve afectado su proyecto personal y profesional, sino también
la capacidad de la institución para acompañar, retener y graduar a sus
estudiantes. Por esta razón, resulta importante desarrollar herramientas
que permitan identificar señales tempranas de riesgo y apoyar procesos
de intervención~\cite{martins2021}.

En este proyecto se propone un sistema de Inteligencia Artificial para
predecir el riesgo de abandono estudiantil a partir de información
académica, económica y personal. El problema se aborda como una tarea de
clasificación binaria, donde el modelo debe determinar si un estudiante
pertenece a la clase \textit{Abandona} o \textit{No abandona}.

La \textbf{pregunta de investigación} planteada es:
\begin{quote}
  \textit{``¿Puede un modelo de Machine Learning predecir el riesgo de
  abandono estudiantil con un F1-score superior al baseline y generar
  recomendaciones personalizadas que apoyen procesos de intervención
  temprana?''}
\end{quote}

Para responder esta pregunta se utilizó el dataset \textit{Predict
Students' Dropout and Academic Success}~\cite{dataset2022}. Se realizó
un análisis exploratorio, entrenamiento de modelos de clasificación,
comparación de métricas y desarrollo de una aplicación en Streamlit.
El documento se organiza así: la Sección~2 describe los datos y el EDA,
la Sección~3 presenta la arquitectura, la Sección~4 detalla la metodología,
la Sección~5 reporta resultados y la Sección~6 discute hallazgos.

%% ==========================================================
%%  2. DATOS Y EDA
%% ==========================================================
\section{Datos y Exploración (EDA)}

\subsection{Descripción del dataset}

El dataset utilizado fue \textit{Predict Students' Dropout and Academic
Success}, el cual contiene información de estudiantes de educación superior
en Portugal. La variable objetivo original incluye tres clases:
\textit{Graduate}, \textit{Dropout} y \textit{Enrolled}. Para este proyecto
se transformó en una variable binaria \texttt{Dropout\_Binary}, donde 1
representa estudiantes que abandonaron y 0 representa graduados o aún
matriculados.

\begin{table}[H]
\centering
\caption{Resumen del dataset utilizado}
\label{tab:dataset}
\begin{tabularx}{\columnwidth}{lX}
\toprule
\textbf{Atributo} & \textbf{Descripción} \\
\midrule
Fuente & UCI Machine Learning Repository \\
Registros (N) & 4\,424 estudiantes \\
Features (p) & 36 variables académicas, económicas y personales \\
Variable objetivo & \texttt{Dropout\_Binary}: 0 = No abandona, 1 = Abandona \\
Distribución & 67.88\% No abandona / 32.12\% Abandona \\
Valores nulos & Ninguno \\
\bottomrule
\end{tabularx}
\end{table}

Después de la transformación, se obtuvo una distribución de 3\,003
estudiantes en la clase \textit{No abandona} y 1\,421 en la clase
\textit{Abandona}.

\subsection{Hallazgos del análisis exploratorio}

El análisis exploratorio mostró que las variables académicas son las más
relacionadas con el abandono estudiantil. Las variables asociadas al número
de unidades curriculares aprobadas y a las calificaciones promedio del
primero y segundo semestre presentaron diferencias importantes entre ambas
clases.

\begin{itemize}[leftmargin=*,topsep=2pt,itemsep=1pt]
  \item \textbf{Variables académicas:} el promedio de unidades aprobadas
        en el 2.º semestre fue 5.61 para quienes no abandonan y 1.94
        para quienes abandonan. La nota promedio del 2.º semestre fue
        12.27 vs. 5.89, respectivamente.
  \item \textbf{Variables económicas:} los estudiantes con matrícula al
        día y sin deudas presentaron menor probabilidad de abandono.
  \item \textbf{Edad de ingreso:} estudiantes que ingresan con mayor
        edad mostraron mayor riesgo, posiblemente por responsabilidades
        laborales o familiares.
  \item \textbf{Sin valores nulos:} no se requirió ninguna estrategia
        de imputación.
\end{itemize}

\begin{figure}[H]
  \centering
  %% Exportar desde el notebook con:
  %% plt.savefig("fig_eda.png", dpi=200, bbox_inches='tight')
  %% y subir a Overleaf con el botón Upload
  \fbox{\rule{0pt}{2.8cm}\rule{\columnwidth-2\fboxsep}{0pt}}
  \caption{Distribución de la variable objetivo y variables más
           correlacionadas con el abandono estudiantil.}
  \label{fig:eda}
\end{figure}

%% ==========================================================
%%  3. ARQUITECTURA DEL SISTEMA
%% ==========================================================
\section{Arquitectura del Sistema}

El sistema está compuesto por tres capas que operan en secuencia:
(1)~pipeline de ML para predicción del riesgo de abandono,
(2)~módulo de recomendación personalizada con estructura de prompt LLM, y
(3)~aplicación web Streamlit para interacción del usuario final
(orientadores y coordinadores académicos).

\begin{figure}[H]
  \centering
  \fbox{\rule{0pt}{3.2cm}\rule{\columnwidth-2\fboxsep}{0pt}}
  \caption{Arquitectura general del sistema. El dataset alimenta el
           pipeline ML; la predicción activa el módulo de recomendación;
           la app Streamlit expone el sistema al usuario.}
  \label{fig:arq}
\end{figure}

\subsection{Componentes principales}

\paragraph{Preprocesamiento.}
Se creó la variable \texttt{Dropout\_Binary} a partir de la variable
original multiclase. El dataset no requirió imputación de valores nulos.
División: 80\% entrenamiento / 20\% prueba con \texttt{random\_state=42}.

\paragraph{Modelo de ML.}
Se evaluaron tres modelos: Baseline (clase mayoritaria), Regresión
Logística (con \texttt{StandardScaler} y \texttt{max\_iter=1000}) y
Random Forest (300 árboles, \texttt{class\_weight=balanced},
\texttt{random\_state=42}). Random Forest fue seleccionado como modelo
principal por su mayor AUC-ROC (0.9316) y su capacidad de manejar el
desbalance de clases.

\paragraph{Componente LLM.}
El módulo de recomendación extrae las variables clave del estudiante
(notas, materias aprobadas, estado financiero, edad de ingreso) y genera
una recomendación de intervención personalizada en lenguaje natural,
clasificando el riesgo en tres niveles: Alto ($\geq 70\%$), Medio
($\geq 40\%$) y Bajo ($< 40\%$). El prompt está estructurado para ser
compatible con LLMs externos (Groq, OpenAI) como extensión futura.

\paragraph{Aplicación Web (Streamlit).}
La interfaz \texttt{app.py} permite al usuario seleccionar un estudiante
por índice, visualizar sus datos, ejecutar la predicción en tiempo real
y obtener la recomendación de intervención. Consume los modelos
serializados con \texttt{joblib}.

%% ==========================================================
%%  4. METODOLOGÍA
%% ==========================================================
\section{Metodología}

\subsection{Configuración experimental}

\begin{table}[H]
\centering
\caption{Configuración de los modelos entrenados}
\label{tab:hparams}
\begin{tabularx}{\columnwidth}{lX}
\toprule
\textbf{Parámetro} & \textbf{Valor} \\
\midrule
Framework & scikit-learn \\
División datos & 80\% train / 20\% test \\
Baseline & DummyClassifier (clase mayoritaria) \\
Regresión Logística & StandardScaler + LR, max\_iter=1000 \\
Random Forest & 300 árboles, class\_weight=balanced \\
Métrica principal & AUC-ROC \\
Hardware & Google Colab (CPU) \\
\bottomrule
\end{tabularx}
\end{table}

\subsection{Estrategia de validación}

Se utilizó \textbf{hold-out} fijo (80/20) dado que el dataset cuenta con
más de 4\,000 registros y la proporción de la clase positiva (32.12\%) es
suficiente para obtener representación en el conjunto de prueba. El uso de
\texttt{class\_weight=balanced} en Random Forest compensa el desbalance
sin necesidad de oversampling explícito.

\subsection{Experimentos realizados}

\begin{enumerate}[leftmargin=*,topsep=2pt,itemsep=2pt]
  \item \textbf{Baseline:} DummyClassifier que predice siempre la clase
        mayoritaria. AUC-ROC = 0.500, F1 = 0.000.
  \item \textbf{Experimento 1 — Regresión Logística:} modelo lineal con
        normalización. AUC-ROC = 0.9266, F1 = 0.8054, Accuracy = 0.8859.
  \item \textbf{Experimento 2 — Random Forest (mejor):} ensamble de 300
        árboles con pesos balanceados. AUC-ROC = 0.9316, F1 = 0.8046,
        Accuracy = 0.8847.
\end{enumerate}

%% ==========================================================
%%  5. RESULTADOS
%% ==========================================================
\section{Resultados}

\subsection{Métricas de evaluación}

\begin{table}[H]
\centering
\caption{Comparación de modelos en el test set (N\,=\,885)}
\label{tab:resultados}
\begin{tabular}{lcccc}
\toprule
\textbf{Modelo} & \textbf{Acc.} & \textbf{F1} & \textbf{AUC} & \textbf{Recall} \\
\midrule
Baseline         & 0.679 & 0.000 & 0.500 & 0.000 \\
Reg. Logística   & 0.886 & 0.805 & 0.927 & 0.736 \\
\textbf{Rand. Forest} & \textbf{0.885} & \textbf{0.805} & \textbf{0.932} & \textbf{0.739} \\
\bottomrule
\end{tabular}
\end{table}

El modelo Random Forest supera al baseline en 43.2 puntos de AUC-ROC y
logra un F1 de 0.805 sobre la clase positiva (abandono). La Regresión
Logística obtiene resultados similares (AUC\,=\,0.927), lo que indica
que las variables del dataset tienen alta separabilidad lineal.
La Tabla~\ref{tab:resultados} confirma que el sistema supera ampliamente
el baseline propuesto.

\subsection{Análisis de resultados}

\begin{figure}[H]
  \centering
  \fbox{\rule{0pt}{2.8cm}\rule{\columnwidth-2\fboxsep}{0pt}}
  \caption{Matriz de confusión del Random Forest. 210 verdaderos positivos
           (abandonos detectados) y 74 falsos negativos.}
  \label{fig:cm}
\end{figure}

\begin{figure}[H]
  \centering
  \fbox{\rule{0pt}{2.5cm}\rule{\columnwidth-2\fboxsep}{0pt}}
  \caption{Curva ROC de los tres modelos. Random Forest alcanza
           AUC\,=\,0.932.}
  \label{fig:roc}
\end{figure}

Las variables con mayor importancia en el Random Forest fueron:
unidades curriculares aprobadas en el 2.º semestre, nota promedio del
2.º semestre, unidades aprobadas en el 1.º semestre, nota promedio del
1.º semestre y estado de pago de matrícula. Estas cinco variables
explican más del 60\% de la importancia total del modelo.

\subsection{Sistema de recomendación}

El módulo de recomendación clasifica el riesgo y genera orientaciones
accionables para cada estudiante:

\begin{lstlisting}[caption={Ejemplo de output del sistema para un
estudiante de alto riesgo}]
Nivel de riesgo: Alto
Probabilidad de abandono: 87.0%

Recomendaciones:
- Asignar tutorias academicas de refuerzo
- Remitir a bienestar universitario
  o apoyo financiero
- Programar revision periodica del caso
\end{lstlisting}

%% ==========================================================
%%  6. DISCUSIÓN
%% ==========================================================
\section{Discusión}

\subsection{Interpretación de resultados}

El sistema logró el objetivo planteado: Random Forest alcanza
AUC-ROC\,=\,0.932, superando el baseline (0.500) en 43 puntos y
respondiendo afirmativamente la pregunta de investigación. La diferencia
marginal con Regresión Logística (0.927) sugiere que el problema tiene
una estructura predominantemente lineal, resultado consistente con la
literatura sobre predicción de abandono~\cite{martins2021}. El módulo
de recomendación cumple su función de traducir la predicción a lenguaje
comprensible para orientadores y coordinadores académicos.

\subsection{Limitaciones}

\begin{itemize}[leftmargin=*,topsep=2pt,itemsep=1pt]
  \item \textbf{LLM simulado:} el componente de recomendación usa reglas
        \texttt{if/else} con estructura de prompt, no un LLM real
        conectado a una API externa. Esto limita la riqueza y
        adaptabilidad de las respuestas.
  \item \textbf{Dominio específico:} el modelo fue entrenado con datos
        de universidades portuguesas; su generalización a contextos
        colombianos requiere re-entrenamiento con datos locales.
  \item \textbf{Sin RAG:} no se implementó recuperación sobre reglamentos
        institucionales, lo que enriquecería las recomendaciones con
        políticas específicas de la universidad.
  \item \textbf{74 falsos negativos:} estudiantes en riesgo que el modelo
        no detecta representan el principal punto de mejora del sistema.
\end{itemize}

\subsection{Trabajo futuro}

Dos extensiones concretas de alto impacto serían: (1)~conectar el
módulo de recomendación a un \textbf{LLM real} vía Groq o OpenAI para
generar orientaciones más ricas y contextualizadas; y (2)~implementar
un módulo \textbf{RAG} que indexe los reglamentos académicos de EAFIT
en ChromaDB, permitiendo que las recomendaciones citen artículos
específicos del reglamento estudiantil.

%% ==========================================================
%%  7. CONCLUSIONES
%% ==========================================================
\section{Conclusiones}

Este proyecto construyó un sistema funcional de predicción de abandono
estudiantil que combina Machine Learning con un módulo de recomendación
en lenguaje natural y una aplicación web desplegable. El modelo Random
Forest alcanzó AUC-ROC\,=\,0.9316 y F1\,=\,0.8046, superando ampliamente
al baseline (AUC\,=\,0.500) y respondiendo afirmativamente la pregunta
de investigación planteada. La aplicación Streamlit permite a orientadores
seleccionar estudiantes y obtener predicciones de riesgo con
recomendaciones de intervención en tiempo real. El principal trabajo futuro
es reemplazar el módulo de recomendación basado en reglas por un LLM real
y añadir RAG sobre reglamentos institucionales para enriquecer las
orientaciones generadas.

%% ==========================================================
%%  CONTRIBUCIONES
%% ==========================================================
\section*{Contribuciones del equipo}

\begin{table}[H]
\centering
\begin{tabular}{lp{5.2cm}}
\toprule
\textbf{Integrante} & \textbf{Contribución principal} \\
\midrule
Santiago Arboleda & EDA, preprocesamiento, visualizaciones \\
Juan Esteban Villada & Modelos ML, evaluación, métricas \\
Cristian Cabezas & Módulo LLM, app Streamlit, integración \\
\bottomrule
\end{tabular}
\caption*{Distribución de responsabilidades por integrante.}
\end{table}

%% ==========================================================
%%  REPOSITORIO Y DEMO
%% ==========================================================
\section*{Repositorio y Demo}

\begin{itemize}[leftmargin=*,topsep=2pt,itemsep=1pt]
  \item \textbf{GitHub:} \url{https://github.com/sarboledag/PROYECTO-FINAL}
  \item \textbf{Video demo ($\leq$3 min):} \url{https://github.com/sarboledag/PROYECTO-FINAL/blob/main/video/video_demo.webm}
  \item \textbf{Instalación:} \texttt{pip install -r requirements.txt}
  \item \textbf{App:} \texttt{streamlit run app.py}
  \item \textbf{Notebook:} \texttt{notebooks/Proyecto\_Final\_IA\_Abandono\_Estudiantil.ipynb}
\end{itemize}

%% ==========================================================
%%  REFERENCIAS
%% ==========================================================
\begin{thebibliography}{99}

\bibitem{dataset2022}
Realinho, V., Vieira Martins, M., Machado, J., \& Baptista, L. (2022).
\textit{Predict Students' Dropout and Academic Success}.
UCI Machine Learning Repository.
\url{https://doi.org/10.24432/C5MC89}

\bibitem{martins2021}
Martins, M. V., Tolledo, D., Machado, J., Baptista, L. M., \& Realinho, V. (2021).
\textit{Early prediction of student's performance in higher education:
A case study}.
Trends and Applications in Information Systems and Technologies, 166--175.

\bibitem{breiman2001}
Breiman, L. (2001).
\textit{Random Forests}.
Machine Learning, 45(1), 5--32.
\url{https://doi.org/10.1023/A:1010933404324}

\bibitem{sklearn2011}
Pedregosa, F., et al. (2011).
\textit{Scikit-learn: Machine Learning in Python}.
Journal of Machine Learning Research, 12, 2825--2830.

\bibitem{streamlit2019}
Streamlit Inc. (2019).
\textit{Streamlit: The fastest way to build data apps}.
\url{https://streamlit.io}

\end{thebibliography}

%% ==========================================================
%%  APÉNDICE
%% ==========================================================
\appendix

\section{Variables más importantes del modelo}

Las 5 variables con mayor importancia según el Random Forest fueron:
(1)~unidades curriculares aprobadas en el 2.º semestre,
(2)~nota promedio del 2.º semestre,
(3)~unidades curriculares aprobadas en el 1.º semestre,
(4)~nota promedio del 1.º semestre, y
(5)~estado de pago de matrícula al día.
Estas variables explican más del 60\% de la importancia total del modelo.

\section{Estructura del repositorio}

\begin{lstlisting}[language=bash,
  caption={Árbol de archivos del proyecto}]
proyecto_final_ia/
|-- app.py
|-- README.md
|-- requirements.txt
|-- data/
|   |-- student_dropout_procesado.csv
|-- models/
|   |-- modelo_random_forest.pkl
|   |-- columnas_modelo.pkl
|-- notebooks/
|   |-- Proyecto_Final_IA_Abandono.ipynb
|-- docs/
|   |-- informe_final.tex
|-- video/
    |-- video_demo.webm
\end{lstlisting}

\end{document}
%% ============================================================
%%  FIN — Universidad EAFIT · IA 2026-1
%% ============================================================